% !TeX program = xelatex
% UTF-8; compile twice with XeLaTeX. No external images or data files needed.
\documentclass[a4paper,landscape,10pt]{article}
\usepackage[margin=14mm,headheight=14pt,headsep=5mm,footskip=7mm]{geometry}
\usepackage{fontspec}
\setmainfont{DejaVu Sans}
\newfontfamily\ipafont{DejaVu Sans}
\usepackage[table]{xcolor}
\usepackage{array,longtable,booktabs,fancyhdr}
\definecolor{ink}{HTML}{18384D}
\definecolor{band}{HTML}{E7EFF4}
\pagestyle{fancy}\fancyhf{}
\fancyhead[L]{\small\color{ink}Parlez Vous QSO?! · Svenska / Swedish}
\fancyhead[R]{\small Englische Phrase · Svensk översättning · IPA}
\fancyfoot[L]{\scriptsize Vorlage: Swedish v 1.0 · Credits: Ullmar -- SM5-1252}
\fancyfoot[R]{\small\thepage}
\setlength{\parindent}{0pt}
\setlength{\tabcolsep}{3mm}
\renewcommand{\arraystretch}{1.25}
\newcommand{\ipa}[1]{{\ipafont[#1]}}
\begin{document}
{\LARGE\bfseries\color{ink}Parlez Vous QSO?!}\hfill{\large Svenska / Swedish}

\vspace{2mm}
{\small QSO-Sprachführer mit schwedischer IPA-Lautschrift. Die Lautschrift bezieht sich auf die schwedische Spalte.}

\vspace{2mm}
{\footnotesize \textbf{Leseschlüssel:} \ipa{ˈ} Hauptbetonung, \ipa{ˌ} Nebenbetonung; \ipa{ː} langer Laut;
\ipa{ɕ} tj-Laut, \ipa{ɧ} sj-Laut; \ipa{ʉː, ɵ} schwedische u-Laute.
Grundlage: mittelschwedisches Standardschwedisch. Die Tonakzente 1 und 2 sind nicht gesondert markiert.}

\vspace{1mm}
\fontsize{9}{11.8}\selectfont
\begin{longtable}{@{}>{\raggedright\arraybackslash}p{13mm}>{\raggedright\arraybackslash}p{73mm}>{\raggedright\arraybackslash}p{75mm}>{\raggedright\arraybackslash}p{90mm}@{}}
\toprule
\textbf{Ref.} & \textbf{English phrase} & \textbf{Svenska} & \textbf{Lautschrift (IPA)}\\
\midrule
\endfirsthead
\toprule
\textbf{Ref.} & \textbf{English phrase} & \textbf{Svenska} & \textbf{Lautschrift (IPA)}\\
\midrule
\endhead
\midrule
\multicolumn{4}{r}{\scriptsize Fortsetzung auf der nächsten Seite}\\
\endfoot
\bottomrule
\endlastfoot
\rowcolor{band} 1.0 & \textbf{Opening Greetings} & \multicolumn{2}{l}{\textbf{Inledande hälsningar}} \\*
1.1 & I am very pleased to meet you. & Jag är mycket glad att träffa dig. & \ipa{jɑːɡ ɛːr ˈmʏkːɛt ɡlɑːd atː ˈtrɛfːa dɛj} \\
1.2 & Thank you for your call & Tack för ditt anrop & \ipa{takː føːr dɪtː ˈanˌruːp} \\
1.3 & Good morning & God morgon & \ipa{ɡuːd ˈmɔrːɔn} \\
1.31 & Good day & Goddag & \ipa{ɡʊˈdɑːɡ} \\
1.32 & Good afternoon & Goddag & \ipa{ɡʊˈdɑːɡ} \\
1.33 & Good evening & God kväll & \ipa{ɡuːd kvɛlː} \\
1.4 & I know only a little Swedish but I enjoy QSOs in Swedish. & Jag kan bara lite svenska men jag gillar att ha QSO på svenska & \ipa{jɑːɡ kanː ˈbɑːra ˈliːtɛ ˈsvɛnːska mɛnː jɑːɡ ˈjɪlːar atː hɑː kʉː ɛs uː poː ˈsvɛnːska} \\
\rowcolor{band} 2.0 & \textbf{Operator} & \multicolumn{2}{l}{\textbf{Operatör}} \\*
2.1 & My name is Peter. & Jag heter Peter & \ipa{jɑːɡ ˈheːtɛr ˈpeːtɛr} \\
2.2 & I have been licensed since 1960. & Jag har haft min licens sedan 1960 & \ipa{jɑːɡ hɑːr haft mɪnː lɪˈsɛnːs ˈseːdan ˌnɪtːɔnˌhɵnːdraˈsɛkːstɪɔ} \\
2.3 & I have been licensed for 43 years. & Jag har varit sändareamatör i 43 år & \ipa{jɑːɡ hɑːr ˈvɑːrɪt ˈsɛnːdarɛamaˌtøːr iː ˌfœrːtɪɔˈtreː oːr} \\
2.4 & I live in Sittingbourne. & Jag bor i Sittingbourne. & \ipa{jɑːɡ buːr iː ˈsɪtɪŋbɔːn} \\
2.5 & Sittingbourne is a village. & Sittingbourne är en by & \ipa{ˈsɪtɪŋbɔːn ɛːr ɛnː byː} \\
2.51 & Sittingbourne is a town. & Sittingbourne är en stad & \ipa{ˈsɪtɪŋbɔːn ɛːr ɛnː stɑːd} \\
2.52 & Sittingbourne is a city. & Sittingbourne är en storstad & \ipa{ˈsɪtɪŋbɔːn ɛːr ɛnː ˈstuːrˌstɑːd} \\
2.6 & Sittingbourne is between London and Dover. & Sittingbourne ligger mellan London och Dover & \ipa{ˈsɪtɪŋbɔːn ˈlɪɡːɛr ˈmɛlːan ˈlɔnːdɔn ɔk ˈdəʊvə} \\
2.7 & I live in the South of England. & Jag bor i södra England & \ipa{jɑːɡ buːr iː ˈsøːdra ˈɛŋːland} \\
2.71 & I live in the North of England. & Jag bor i norra England & \ipa{jɑːɡ buːr iː ˈnɔrːa ˈɛŋːland} \\
2.72 & I live in the East of England. & Jag bor i östra England & \ipa{jɑːɡ buːr iː ˈœsːtra ˈɛŋːland} \\
2.73 & I live in West of England. & Jag bor i västra England. & \ipa{jɑːɡ buːr iː ˈvɛsːtra ˈɛŋːland} \\
2.74 & I live in South East England. & Jag bor i sydöstra England & \ipa{jɑːɡ buːr iː ˈsyːdˌœsːtra ˈɛŋːland} \\
2.75 & I live in South West England. & Jag bor i sydvästra England. & \ipa{jɑːɡ buːr iː ˈsyːdˌvɛsːtra ˈɛŋːland} \\
2.76 & I live in North East England. & Jag bor i nordöstra England & \ipa{jɑːɡ buːr iː ˈnuːrdˌœsːtra ˈɛŋːland} \\
2.77 & I live in North West England. & Jag bor i nordvästra England & \ipa{jɑːɡ buːr iː ˈnuːrdˌvɛsːtra ˈɛŋːland} \\
2.78 & I live in central England. & Jag bor i mitten av England & \ipa{jɑːɡ buːr iː ˈmɪtːɛn ɑːv ˈɛŋːland} \\
2.8 & I live in the country. & Jag bor på landet & \ipa{jɑːɡ buːr poː ˈlanːdɛt} \\
2.81 & I live by the sea. & Jag bor vid havet. & \ipa{jɑːɡ buːr viːd ˈhɑːvɛt} \\
2.82 & I live near mountains. & Jag bor nära bergen & \ipa{jɑːɡ buːr ˈnɛːra ˈbærːjɛn} \\
2.9 & Sittingbourne is 50Km from London. & Sittingbourne ligger 50 km från London & \ipa{ˈsɪtɪŋbɔːn ˈlɪɡːɛr ˈfɛmːtɪɔ kɪlɔˈmeːtɛr froːn ˈlɔnːdɔn} \\
2.11 & I am retired. & Jag är pensionerad & \ipa{jɑːɡ ɛːr paŋɧʊˈneːrad} \\
2.12 & My home call sign is M0ECS. & Min hemmasignal är M0ECS & \ipa{mɪnː ˈhɛmːasɪŋˌnɑːl ɛːr ɛm ˈnɵlː ɛː seː ɛs} \\
\rowcolor{band} 3.0 & \textbf{Station} & \multicolumn{2}{l}{\textbf{Station}} \\*
3.1 & My radio is an Icom 706. & Min station är en ICOM sjuhundrasex & \ipa{mɪnː staˈɧuːn ɛːr ɛnː ˈaɪkɒm ˌɧʉːhɵnːdraˈsɛkːs} \\
3.2 & My station is all homebrew (homemade). & Min station är helt hembyggd & \ipa{mɪnː staˈɧuːn ɛːr heːlt ˈhɛmːˌbʏɡːd} \\
3.3 & I am running 5 Watts. & Jag kör med 5 watt. & \ipa{jɑːɡ ɕøːr meːd fɛmː vatː} \\
3.4 & My antenna is a long wire. & Min antenn är en longwire & \ipa{mɪnː anˈtɛnː ɛːr ɛnː ˈlɒŋwaɪə} \\
3.41 & My antenna is a dipole. & Min antenn är en dipol & \ipa{mɪnː anˈtɛnː ɛːr ɛnː dɪˈpoːl} \\
3.42 & My antenna is a G5RV. & Min antenn är en G5RV. & \ipa{mɪnː anˈtɛnː ɛːr ɛnː ɡeː fɛmː ɛr veː} \\
3.43 & My antenna is a doublet. & Min antenn är en doublet. & \ipa{mɪnː anˈtɛnː ɛːr ɛnː ˈdʌblət} \\
3.44 & My antenna is a delta loop. & Min antenn är en deltaloop & \ipa{mɪnː anˈtɛnː ɛːr ɛnː ˈdɛlːtaˌluːp} \\
3.45 & My antenna is a vertical. & Min antenn är en vertikal & \ipa{mɪnː anˈtɛnː ɛːr ɛnː vɛrtɪˈkɑːl} \\
3.46 & My antenna is a log periodic. & Min antenn är en log-periodic. & \ipa{mɪnː anˈtɛnː ɛːr ɛnː lɒɡ pɪəɹiˈɒdɪk} \\
3.47 & My antenna is a Yagi beam. & Min antenn är en yagi & \ipa{mɪnː anˈtɛnː ɛːr ɛnː ˈjɑːɡɪ} \\
3.48 & My antenna is a 5 element beam. & Min antenn är en 5 elements beam & \ipa{mɪnː anˈtɛnː ɛːr ɛnː fɛmː ɛlɛˈmɛnːts biːm} \\
\rowcolor{band} 4.0 & \textbf{Conditions} & \multicolumn{2}{l}{\textbf{Konditionerna}} \\*
4.1 & Thank you for my 599 report. & Tack för min rapport 599 & \ipa{takː føːr mɪnː raˈpɔrːt fɛmː ˈniːɔ ˈniːɔ} \\
4.2 & The band conditions are good. & Det är bra konditioner & \ipa{deː ɛːr brɑː kɔndɪˈɧuːnɛr} \\
4.3 & The band conditions are bad. & Konditionerna på bandet är dåliga & \ipa{kɔndɪˈɧuːnɛrna poː ˈbanːdɛt ɛːr ˈdoːlɪɡa} \\
4.4 & There is strong interference here. & Jag har starka störningar här & \ipa{jɑːɡ hɑːr ˈstarːka ˈstøːrˌnɪŋːar hɛːr} \\
4.5 & Can you change to the 80m band? & Kan du QSY-a till 80 m? & \ipa{kanː dʉː kʉː ɛs yː a tɪlː ˈɔtːɪɔ ˈmeːtɛr} \\
4.6 & Please change frequency to 3.558 MHz. & Var snäll och QSY-a till 3,558 MHz. & \ipa{vɑːr snɛlː ɔk kʉː ɛs yː a tɪlː treː ˈkɔmːa fɛmː fɛmː ˈɔtːa ˈmeːɡaˌhærːts} \\
\rowcolor{band} 5.0 & \textbf{General Queries and Comments} & \multicolumn{2}{l}{\textbf{Allmänna frågor och kommentarer}} \\*
5.1 & What does <\#\#\#\#\#> mean please? & Vad betyder <\#\#\#\#\#>? & \ipa{vɑːd bɛˈtyːdɛr …} \\
5.2 & Please wait a moment, I need to check my Swedish dictionary. & Vänta lite, jag måste kolla min svenska ordbok & \ipa{ˈvɛnːta ˈliːtɛ jɑːɡ ˈmɔsːtɛ ˈkɔlːa mɪnː ˈsvɛnːska ˈuːrdˌbuːk} \\
5.3 & How do you say, <\#\#\#\#\#> in Swedish ? & Hur säger man <\#\#\#\#\#> på svenska? & \ipa{hʉːr ˈsɛjːɛr manː … poː ˈsvɛnːska} \\
5.5 & My favourite modes of operation are CW / RTTY / PSK / MT63 / Olivia / Hellschreiber / Digital Modes / and SSB & Jag tycker mest om att köra CW / RTTY / PSK / MT63 / Olivia / Hellschreiber / digitala moder /  och SSB & \ipa{jɑːɡ ˈtʏkːɛr mɛst ɔmː atː ˈɕøːra seː ˈdɵbːɛl veː / ɛr teː teː yː / peː ɛs koː / ɛm teː ˌsɛkːstɪɔˈtreː / ɔˈliːvɪa / ˈhɛlˌʃʁaɪbɐ / dɪɡɪˈtɑːla ˈmuːdɛr / ɔk ɛs ɛs beː} \\
5.6 & My best DX is Australia. & Mitt bästa DX är Australien & \ipa{mɪtː ˈbɛsːta deː ɛks ɛːr aʊˈstrɑːlɪɛn} \\
5.7 & I enjoy working contests. & Jag gillar att köra contest & \ipa{jɑːɡ ˈjɪlːar atː ˈɕøːra ˈkɒntɛst} \\
5.8 & I like working QRP. & Jag tycker om att köra QRP & \ipa{jɑːɡ ˈtʏkːɛr ɔmː atː ˈɕøːra kʉː ɛr peː} \\
\rowcolor{band} 6.0 & \textbf{The Weather} & \multicolumn{2}{l}{\textbf{Vädret}} \\*
6.1 & It is sunny. & Det är soligt. & \ipa{deː ɛːr ˈsuːlɪt} \\
6.2 & It is windy. & Det är blåsigt & \ipa{deː ɛːr ˈbloːsɪt} \\
6.3 & It is raining. & Det regnar & \ipa{deː ˈrɛŋːnar} \\
6.4 & It is snowing. & Det snöar. & \ipa{deː ˈsnøːar} \\
6.5 & It is very hot. & Det är väldigt varmt & \ipa{deː ɛːr ˈvɛlːdɪt varːmt} \\
6.6 & It is very cold. & Det är väldigt kallt & \ipa{deː ɛːr ˈvɛlːdɪt kalːt} \\
6.7 & The temperature is 8C. & Temperaturen är åtta grader celsius & \ipa{tɛmpɛraˈtʉːrɛn ɛːr ˈɔtːa ˈɡrɑːdɛr ˈsɛlːsjɵs} \\
6.8 & It is overcast. & Det är mulet & \ipa{deː ɛːr ˈmʉːlɛt} \\
6.9 & It is foggy. & Det är dimmigt & \ipa{deː ɛːr ˈdɪmːɪt} \\
6.11 & It has been raining all day. & Det har regnat hela dagen & \ipa{deː hɑːr ˈrɛŋːnat ˈheːla ˈdɑːɡɛn} \\
6.12 & There is 20cm of snow here. & Vi har 20 cm snö här & \ipa{viː hɑːr ˈɕʉːɡɔ sɛntɪˈmeːtɛr snøː hɛːr} \\
\rowcolor{band} 8.0 & \textbf{QSL} & \multicolumn{2}{l}{\textbf{QSL}} \\*
8.1 & My QSL card is sure via the bureau & Mitt QSL-kort kommer via byrån. & \ipa{mɪtː kʉː ɛs ɛl kɔrːt ˈkɔmːɛr ˈviːa ˈbyːroːn} \\
8.2 & QSL via eQSL & Jag skickar ett eQSL & \ipa{jɑːɡ ˈɧɪkːar ɛtː iː kʉː ɛs ɛl} \\
8.3 & No eQSL please & Inget eQSL, tack! & \ipa{ˈɪŋːɛt iː kʉː ɛs ɛl takː} \\
8.4 & I will send you my card direct. & Jag skickar dig mitt kort direkt. & \ipa{jɑːɡ ˈɧɪkːar dɛj mɪtː kɔrːt dɪˈrɛkːt} \\
8.5 & What is your address? & Vad är din adress? & \ipa{vɑːd ɛːr dɪnː aˈdrɛsː} \\
8.6 & Please send your QSL card direct. & Var snäll och skicka mig ditt QSL direkt & \ipa{vɑːr snɛlː ɔk ˈɧɪkːa mɛj dɪtː kʉː ɛs ɛl dɪˈrɛkːt} \\
8.7 & My address is... & Min adress är… & \ipa{mɪnː aˈdrɛsː ɛːr …} \\
8.8 & My address is shown on QSL.com / EHam.net & Min adress finns på QSL.com / eHam.net. & \ipa{mɪnː aˈdrɛsː fɪnːs poː kʉː ɛs ɛl pɵŋkt kɔm / iː hæm pɵŋkt nɛt} \\
8.9 & Sorry I do not send QSL cards. & Jag är ledsen men jag skickar inga QSL & \ipa{jɑːɡ ɛːr ˈlɛsːɛn mɛnː jɑːɡ ˈɧɪkːar ˈɪŋːa kʉː ɛs ɛl} \\
8.11 & Please send QSL card via G1JQH. & Var snäll och skicka QSL via G1JQH. & \ipa{vɑːr snɛlː ɔk ˈɧɪkːa kʉː ɛs ɛl ˈviːa ɡeː ɛtː jiː kʉː hoː} \\
8.12 & My WACRAL / GQRP / DXCC / FISTS number is 1234. & Mitt WACRAL / GQRP / DXCC / FISTS nummer är 1234. & \ipa{mɪtː ˈdɵbːɛl veː ɑː seː ɛr ɛl / ɡeː kʉː ɛr peː / deː ɛks seː seː / fɪsts ˈnɵmːɛr ɛːr ˌɛtːˈtʉːsɛn ˌtvoːˈhɵnːdra ˌtrɛtːɪɔˈfyːra} \\
\rowcolor{band} 9.0 & \textbf{Final Greetings} & \multicolumn{2}{l}{\textbf{Avskedshälsningar}} \\*
9.1 & Thank you. & Tack så mycket & \ipa{takː soː ˈmʏkːɛt} \\
9.2 & Thank you for a very enjoyable contact. & Tack för en mycket trevlig kontakt & \ipa{takː føːr ɛnː ˈmʏkːɛt ˈtreːvlɪɡ kɔnˈtakːt} \\
9.3 & I enjoyed our chat. & Jag har verkligen uppskattat vår pratstund & \ipa{jɑːɡ hɑːr ˈværːklɪɡɛn ˈɵpːˌskatːat voːr ˈprɑːtˌstɵnːd} \\
9.4 & I enjoyed our QSO in Swedish. & Jag uppskattade vårt QSO på svenska & \ipa{jɑːɡ ˈɵpːˌskatːadɛ voːrt kʉː ɛs uː poː ˈsvɛnːska} \\
9.5 & It has been great fun. & Det var väldigt roligt.\newline Det var riktigt kul. (freier) & \ipa{deː vɑːr ˈvɛlːdɪt ˈruːlɪt / deː vɑːr ˈrɪkːtɪt kʉːl} \\
9.6 & If you would like a regular contact with me, you can contact me via email. & Om du vill ha kontakt med mig regelbundet kan du kontakta mig via e-post. & \ipa{ɔmː dʉː vɪlː hɑː kɔnˈtakːt meːd mɛj ˈreːɡɛlˌbɵnːdɛt kanː dʉː kɔnˈtakːta mɛj ˈviːa ˈiːˌpɔsːt} \\
9.7 & My email address is... & Min e-postadress är…. & \ipa{mɪnː ˈiːpɔstˌadrɛsː ɛːr …} \\
9.8 & My email address is shown on QSL.com / EHam.net & Min e-postadress finns på QSL.com / EHam.net & \ipa{mɪnː ˈiːpɔstˌadrɛsː fɪnːs poː kʉː ɛs ɛl pɵŋkt kɔm / iː hæm pɵŋkt nɛt} \\
9.9 & I look forward to meeting you again. & Jag hoppas vi träffas igen & \ipa{jɑːɡ ˈhɔpːas viː ˈtrɛfːas ɪˈjɛnː} \\
9.11 & My best wishes to you and your family. & Mina bästa hälsningar till dig och din familj & \ipa{ˈmiːna ˈbɛsːta ˈhɛlːsˌnɪŋːar tɪlː dɛj ɔk dɪnː faˈmɪlːj} \\
9.12 & Best wishes to your wife. & Bästa hälsningar till din fru (XYL) & \ipa{ˈbɛsːta ˈhɛlːsˌnɪŋːar tɪlː dɪnː frʉː (ɛks yː ɛl)} \\
9.121 & Best wishes to your husband. & Bästa hälsningar till din man & \ipa{ˈbɛsːta ˈhɛlːsˌnɪŋːar tɪlː dɪnː manː} \\
9.13 & Good luck. & Lycka till & \ipa{ˈlʏkːa tɪlː} \\
9.14 & Good health. & God hälsa & \ipa{ɡuːd ˈhɛlːsa} \\
9.15 & Goodbye & Adjö / På återhörande / Hej då! (frei) & \ipa{aˈjøː / poː ˈoːtɛrˌhøːrandɛ / hɛj doː} \\
9.16 & I wish you an enjoyable afternoon. & Ha en trevlig eftermiddag & \ipa{hɑː ɛnː ˈtreːvlɪɡ ˈɛfːtɛrˌmɪdːaɡ} \\
9.161 & I wish you an enjoyable evening. & Ha en trevlig kväll & \ipa{hɑː ɛnː ˈtreːvlɪɡ kvɛlː} \\
9.162 & I wish you an enjoyable weekend. & Jag önskar dig en trevlig veckohelg & \ipa{jɑːɡ ˈœnːskar dɛj ɛnː ˈtreːvlɪɡ ˈvɛkːɔˌhɛlːj} \\
9.17 & Enjoy your holiday. & Jag önskar dig en skön semester. & \ipa{jɑːɡ ˈœnːskar dɛj ɛnː ɧøːn sɛˈmɛsːtɛr} \\
9.18 & Happy New Year. & Gott Nytt År & \ipa{ɡɔtː nʏtː oːr} \\
\rowcolor{band} 10.0 & \textbf{Country} & \multicolumn{2}{l}{\textbf{Land}} \\*
10.1 & Albania & Albanien & \ipa{alˈbɑːnɪɛn} \\
10.2 & Austria & Österrike & \ipa{ˈœsːtɛrˌriːkɛ} \\
10.3 & Belgium & Belgien & \ipa{ˈbɛlːɡɪɛn} \\
10.4 & Belorussia & Belarus/Vitryssland & \ipa{bɛlaˈrʉːs / ˈviːtˌrʏsːland} \\
10.5 & Bosnia \& Herzegovina & Bosnien-Herzegovina & \ipa{ˈbɔsːnɪɛn hɛrtsɛɡɔˈviːna} \\
10.6 & Bulgaria & Bulgarien & \ipa{bɵlˈɡɑːrɪɛn} \\
10.7 & Croatia & Kroatien & \ipa{krɔˈɑːtsɪɛn} \\
10.8 & Czech Republic & Tjeckien & \ipa{ˈɕɛkːɪɛn} \\
10.9 & Denmark & Danmark & \ipa{ˈdanːmark} \\
10.11 & Eire & Irland /Eire & \ipa{ˈiːrland / ˈeːɾʲə} \\
10.12 & England & England & \ipa{ˈɛŋːland} \\
10.13 & Estonia & Estland & \ipa{ˈɛstːland} \\
10.14 & Finland & Finland & \ipa{ˈfɪnːland} \\
10.15 & France & Frankrike & \ipa{ˈfraŋːkˌriːkɛ} \\
10.16 & Germany & Tyskland & \ipa{ˈtʏsːkland} \\
10.17 & Gibraltar & Gibraltar & \ipa{jɪˈbraltar} \\
10.18 & Greece & Grekland & \ipa{ˈɡreːkland} \\
10.19 & Holland & Holland & \ipa{ˈhɔlːand} \\
10.21 & Hungary & Ungern & \ipa{ˈɵŋːɛrn} \\
10.22 & Iceland & Island & \ipa{ˈiːsland} \\
10.23 & Northern Ireland & Nord-Irland & \ipa{ˈnuːrdˌiːrland} \\
10.24 & Isle of Man & Isle-of-Man & \ipa{ˌaɪl əv ˈmæn} \\
10.25 & Italy & Italien & \ipa{ɪˈtɑːlɪɛn} \\
10.26 & Latvia & Lettland & \ipa{ˈlɛtːland} \\
10.27 & Liechtenstein & Liechtenstein & \ipa{ˈliçtn̩ʃtaɪn} \\
10.28 & Lithuania & Litauen & \ipa{lɪˈtauɛn} \\
10.29 & Luxembourg & Luxemburg & \ipa{ˈlɵkːsɛmˌbɵrːɡ} \\
10.31 & Macedonia & Makedonien & \ipa{makɛˈduːnɪɛn} \\
10.32 & Moldavia & Moldavien / Moldavia & \ipa{mɔlˈdɑːvɪɛn / mɔlˈdɑːvɪa} \\
10.33 & Netherlands & Nederländerna & \ipa{ˈneːdɛrˌlɛnːdɛrna} \\
10.34 & Norway & Norge & \ipa{ˈnɔrːjɛ} \\
10.35 & Poland & Polen & \ipa{ˈpuːlɛn} \\
10.36 & Portugal & Portugal & \ipa{pɔrtɵˈɡɑːl} \\
10.37 & Romania & Rumänien & \ipa{rɵˈmɛːnɪɛn} \\
10.38 & Russian Federation & Ryssland & \ipa{ˈrʏsːland} \\
10.39 & Scotland & Skottland & \ipa{ˈskɔtːland} \\
10.41 & Slovakia & Slovakien & \ipa{slɔˈvɑːkɪɛn} \\
10.42 & Slovenia & Slovenien & \ipa{slɔˈveːnɪɛn} \\
10.43 & Spain & Spanien & \ipa{ˈspɑːnɪɛn} \\
10.44 & Svalbard & Svalbard & \ipa{ˈsvɑːlbard} \\
10.45 & Sweden & Sverige & \ipa{ˈsværːjɛ} \\
10.46 & Switzerland & Schweiz & \ipa{ʂvɛjts} \\
10.47 & Turkey & Turkiet & \ipa{ˈtɵrːkɪɛt} \\
10.48 & Ukraine & Ukraina & \ipa{ʉːˈkrɑjːna} \\
10.49 & Wales & Wales & \ipa{weɪlz} \\
\end{longtable}
\normalsize
\section*{Hinweise zur Übertragung und Aussprache}
\fontsize{8.5}{10}\selectfont\raggedright
\textbf{Vorlage und Umfang.} Grundlage ist „Swedish v 1.0.doc“; Credit der Quelle: Ullmar, SM5-1252.
Alle 142 Phrasen und Ländernamen einschließlich ihrer Alternativen sind enthalten. Nummernsprünge und der fehlende Abschnitt 7 bleiben erhalten.

\smallskip
\textbf{Dokumentierte Korrekturen.} Referenz „132“ wurde zu „1.32“, „G5GV“ zu „G5RV“
(3.42, entsprechend der englischen Spalte) und „address“ zu „adress“ (8.5) berichtigt.
Mehrfache Leerzeichen wurden vereinheitlicht; die Liste in 5.5 erhielt Leerzeichen um die Schrägstriche.
Der abgesetzte Gedankenstrich vor „nummer“ in 8.12 entfiel. Die englischen Erläuterungen „More freely“
und „freely“ in 9.5 und 9.15 wurden als „freier“ und „frei“ wiedergegeben und die Alternativen deutlicher
getrennt. Diese Erläuterungen werden nicht mitgesprochen.

\smallskip
\textbf{Inhaltliche Eigenheiten.} Die übrigen Formulierungen bleiben wie in der Quelle. Bei 3.1 steht „Min station“ statt „My radio“.
Belarus / Vitryssland, Irland / Eire und Moldavien / Moldavia sind als Alternativen übernommen.
Historische Ländernamen und Webadressen wurden nicht modernisiert. Das in Klammern stehende XYL in
9.12 ist als optional gesprochenes Kürzel transkribiert.

\smallskip
\textbf{Aussprachebasis.} Die IPA-Spalte ist eine breite, manuell erstellte Aussprachehilfe für
mittelschwedisches Standardschwedisch, keine enge Transkription einer Aufnahme. \ipa{ˈ} markiert die
Haupt-, \ipa{ˌ} ausgewählte Nebenbetonungen. \ipa{ː} kann sowohl lange Vokale als auch lange Konsonanten
kennzeichnen. Die u-Laute \ipa{ʉː, ɵ} sind zentraler als deutsches u. Regionale Varianten sind nicht vollständig dargestellt.

\smallskip
\textbf{Wortakzent und Satzmelodie.} Viele schwedische Varianten unterscheiden Akzent 1 und Akzent 2
auch durch den Tonverlauf. Die Zeichen \ipa{ˈ, ˌ} markieren hier ausschließlich Betonung, nicht diese
Tonakzente. Tonverläufe sind nicht kodiert und variieren regional sowie mit dem Satzkontext; zusätzliches Hören
und Nachsprechen hilft bei der natürlichen Aussprache.

\smallskip
\textbf{Besondere Laute und Reduktionen.} \ipa{ɕ} steht für den tj-Laut, etwa in „kör“ und „tjugo“;
\ipa{ɧ} für eine gebräuchliche mittelschwedische sj-Lautvariante, etwa in „station“, „skickar“ und „sjuhundra“.
Bei „Schweiz“ ist eine Variante mit \ipa{ʂ} gewählt. Kombinationen von r mit t, d, n, s oder l können
retroflex gesprochen werden; die Tabelle lässt viele davon als r plus Konsonant stehen.
„mig / dig“ werden als \ipa{mɛj / dɛj} gelesen. „det“ ist als \ipa{deː}, „ledsen“ ohne d und Endungen
wie -igt überwiegend ohne g transkribiert. Andere Wörter sind bewusst in deutlicher Lesesprache dargestellt:
„jag“ mit g, „och“ mit k sowie vollere Zahlenformen mit -tio. Im Gespräch sind weitere Kürzungen üblich.


\smallskip
\textbf{Namen und Kürzel.} Sittingbourne, Dover und die Fachwörter Longwire, Doublet, Log-Periodic,
Beam und Contest stehen mit englischer Aussprache. Icom verwendet ebenfalls eine englische Variante.
Isle of Man und Wales sind englisch, Eire irisch, Hellschreiber und Liechtenstein deutsch transkribiert;
andere Namen sind schwedisch beziehungsweise schwedisch angepasst gelesen. Rufzeichen und die meisten Kürzel stehen mit schwedischen Buchstabennamen,
nicht im internationalen Buchstabieralphabet: Q „ku“, W „dubbel-ve“, J „ji“, H „hå“.
Das e in eQSL, eHam und e-post ist \ipa{iː}; FISTS wird als Wort gesprochen.
QSY-a wird hier als „Q S Y“ mit angehängtem a gelesen.

\smallskip
\textbf{Zahlen und Platzhalter.} Zahlen sind ausgeschrieben: 1960 „nittonhundrasextio“,
43 „fyrtiotre“, 706 entsprechend der Quelle „sjuhundrasex“, 599 „fem nio nio“ und
1234 „ett tusen två hundra trettiofyra“. 3,558 MHz steht für „tre komma fem fem åtta megahertz“.
80 m und 20 cm werden als „åttio meter“ und „tjugo centimeter“ gelesen.
Schrägstriche trennen Alternativen oder Listenelemente. Platzhalter \texttt{<\#\#\#\#\#>}
erscheinen in der Lautschrift als Auslassung und werden nicht mitgesprochen.

\smallskip
\textbf{Fachquellen.} Olle Engstrand: „Swedish“, Journal of the International Phonetic Association
20(1), 1990, S. 42--44, DOI 10.1017/S0025100300004060.
Gösta Bruce, Claes-Christian Elert, Olle Engstrand und Pär Wretling: „Phonetics and Phonology of the
Swedish Dialects“, ICPhS 1999, zur regionalen Realisierung der Wortakzente.
Die einzelnen Phrasen wurden manuell transkribiert.

\end{document}
